\subsection{A conditional extension for observer curves}
\label{sec:observer-families}

The private-handle count can be supplemented by certain exact incidence
data when those data are already available.  The following algebraic
observation identifies a limited interface for future curve transport;
it does not provide the missing initial geometric data.

\begin{lemma}[Observer families under column peeling]
\label{obs:lem-families}
Let $A$ be the nonnegative integer incidence matrix of $m$ task rows,
and let $Y$ be a nonnegative integer matrix of observer rows on the
same columns.  Suppose the private columns are partitioned into
families $F_\ell$ such that all columns in one family are identical
across both $A$ and $Y$ initially.  Assume every private task entry
is zero or at least two; observer rows are never eligible for pivots.

Under any positive-pivot schedule on active columns, equality within
each private family is preserved.  One representative column and
the binary multiplicity $|F_\ell|$ therefore suffice for that family.
If the schedule uses task column leaves, let $W\geq2$ bound every
\emph{exact} initial entry of $A$ and $Y$.  After $t\leq m$ pivots,
each surviving observer entry is at most
\[
 W(1+W)^t,
\]
and hence has $O((m+1)\log(W+1))$ bits, independently of the
private-family multiplicities.
\end{lemma}

\begin{proof}
For a pivot $(r,g)$ the observer update is
\begin{equation}
 Y'_{u,h}=Y_{u,h}+Y_{u,g}A_{r,h}\qquad(h\ne g).
 \label{obs:eq-update}
\end{equation}
The task rows follow the same positive update rule.  Equal columns
receive equal additions in every surviving row, so family equality
is invariant.  Private task columns cannot acquire a unit entry,
by Lemma~\ref{lem:inert-columns}, and their multiplicities are unchanged.

For a task column leaf, $A_{r,g}=1$ and $A_{i,g}=0$ for every other
live task row $i$.  Thus task rows never fill: every surviving task
entry is an initial entry and is at most $W$.  If $M_t$ bounds the
observer entries after $t$ pivots, \eqref{obs:eq-update} gives
$M_{t+1}\leq(1+W)M_t$.  Starting with $M_0\leq W$ proves the bound.
No update sums over the columns of a private family, so its
multiplicity contributes no factor to this recurrence.
\end{proof}

\paragraph{Geometric interpretation.}
Suppose the observer rows are authenticated unsigned cocore
incidences of embedded curves disjoint from all task attaching
curves.  Choose the attaching collars sufficiently small to avoid
the observers.  The geometric positive-pivot proof can likewise
choose its sliding bands in a sufficiently small neighbourhood of
the donor task curve, avoiding the finite union of observer curves.  Each disc
slide then adds one copy of the relevant cocore's observer
intersections, giving exactly \eqref{obs:eq-update}.  The observers
survive the attachment and cancellation because they avoid the
attaching region.  This applies the same embedded disc-slide
argument as the task-row formula~\cite{marked-parallelity-85}.

\paragraph{What additional input is required.}
The update must use \emph{uncapped} donor entries: a value recording
only ``at least two'' suffices for pivot admissibility but not for
exact observer counts.  The number of families and their initial
observer coordinates must also be supplied and geometrically
authenticated.  Scalar genus and multiplicity profiles do not
supply them or bound the number of required families.  Multiplicity
itself still requires $\log(|F_\ell|+1)$ bits of storage.  Even when
the hypotheses hold, the output consists of cocore intersection
counts; it does not determine cyclic order, source addresses, or
the embedded position of an arbitrary boundary pattern.  The
current implementation does not produce these initial observer
data or implement this conditional extension.
